\documentclass[12pt]{article}

\usepackage{graphicx}
\usepackage{amssymb}
\usepackage{amsmath}
\usepackage[margin=1in]{geometry}
\usepackage{fancyhdr}
\usepackage{enumerate}
\usepackage[shortlabels]{enumitem}

\pagestyle{fancy}
\fancyhead[l]{Li Yifeng}
\fancyhead[c]{Homework \#1}
\fancyhead[r]{\today}
\fancyfoot[c]{\thepage}
\renewcommand{\headrulewidth}{0.2pt}
\setlength{\headheight}{15pt}

\begin{document}
	
	\section*{Exercise 1}
	
	\noindent Find an application of propositional logic in computer science or artificial intelligence and describe it in $150$ to $300$ words. Share your answer with your peers on the Moodle discussion forum.
	
	\bigskip
	
	\begin{enumerate}[label={},leftmargin=0in]\item
		\subsection*{Answer}
		
			I think that propositional logic is the base of Logical Programming, useful for languages like Prolog and so on, providing a traceable way of modeling.
			
			Besides, if we look at the area of Computer Architecture, logic gates for example, is nearly identical to propositional logic. A deep understanding is essential for building complex computer systems.
			
			In conclusion, we can say that propositional logic is the computational foundation of Computer Science.
			
	\end{enumerate}
	
\end{document}
